\section{Giới thiệu}
\label{sec:introduction}

Trong quá trình phát triển của ngành Trí tuệ Nhân tạo, các giải thuật
tìm kiếm trong không gian trạng thái đóng vai trò là vô cùng quan trọng
vì nó là nền tảng cốt lõi cho việc giải quyết các bài toán lập kế hoạch
tự động, tối ưu hóa tổ hợp và ra quyết định thông minh
\cite{russell2020aima, edelkamp2011heuristic}.
Để đánh giá một cách rõ ràng và đầy đủ được cho phần hiệu năng của các giải
thuật tìm kiếm không có thông tin và tìm kiếm có thông tin thì các trò chơi
giải đố tổ hợp kinh điển (chẳng hạn như 15-Puzzle, Sokoban, Rubik's Cube)
thường xuyên được sử dụng làm môi trường chuẩn đánh giá
\cite{pearl1984heuristics, sturtevant2012benchmarks}.

Trong số các trò chơi giải đố không gian, \textbf{Bloxorz}, được thiết
kế bởi Damien Clarke vào năm 2007 \cite{clarke2007bloxorz}, nổi lên như
một bài toán mô phỏng hay và đầy thách thức đối với các hệ thống AI.
Trò chơi đặt người chơi vào nhiệm vụ điều khiển một khối hình hộp chữ nhật
đồng chất kích thước $1 \times 1 \times 2$ lăn lật trên một mặt sàn lưới
rời rạc $\mathbb{Z}^2$ để đưa khối về vị trí đích ở tư thế dựng đứng.
Điểm hấp dẫn đặc biệt của Bloxorz nằm ở chỗ: bài toán không đơn thuần là
điều hướng một chất điểm trong mê cung 2D tĩnh, mà là làm sao để 
vật thể cứng có kích thước hình học có thể di chuyển về đích mà không bị
gặp bẫy, các cảm biến hoặc các ô đặc biệt. Ta có các tư thế không gian (đứng,
nằm ngang, nằm dọc) ảnh hưởng trực tiếp từ hướng đi cho tới diện tích tiếp xúc,
khả năng chịu tải của bề mặt sàn và các cơ chế kích hoạt tương tác phức tạp.

\subsection{Thách thức và Độ phức tạp}
Mặc dù luật chơi bề ngoài rất dễ hiểu nhưng bài toán Bloxorz tiềm ẩn độ phức
tạp tính toán rất cao:
\begin{enumerate}[1.]
    \item \textbf{Độ phức tạp PSPACE-Complete:}
    Nghiên cứu lý thuyết độ phức tạp tính toán của van der Zanden và
    Bodlaender (2014) \cite{vanderzanden2014bloxorz} đã chứng minh
    rằng việc xác định tính khả giải của một màn chơi Bloxorz tổng
    quát có chứa các công tắc và cầu lật là một bài toán thuộc lớp
    \textbf{PSPACE-complete} \cite{sipser2013computation, hearn2009games}
     tập hợp các bài toán quyết định khó nhất đòi hỏi lượng bộ nhớ
    không gian đa thức nhưng thời gian thực thi có thể bùng nổ theo cấp
    số mũ. Tương tự các nghiên cứu về câu đố lăn khối đa diện trên lưới
    của Buchin và cộng sự \cite{buchin2007rolling}, không gian trạng
    thái của Bloxorz bùng nổ tổ hợp theo quy mô bản đồ và số lượng cơ
    chế kích hoạt.

    \item \textbf{Môi trường đa Trạng thái và biến động:}
    Với các màn chơi có sự xuất hiện của các ô gạch cam chịu lực kém
    phụ thuộc vào tư chế tiếp xúc mà dẫn tới các hệ thống cầu đóng/mở
    đan xen chịu sự điều khiển của nhiều loại công tắc cơ học (công tắc
    tròn tiếp xúc nhẹ, công tắc chữ X chịu lực tập trung), và đặc biệt
    là công tắc phân thân chia tách khối thành hai khối lập phương
    $1 \times 1 \times 1$ độc lập khiến không gian trạng thái mở rộng
    thành một tích Descartes phức hợp giữa cấu hình hình học của khối
    và véc-tơ trạng thái của toàn bộ môi trường.

    \item \textbf{Thách thức trong Thiết kế Hàm Heuristic:}
    Để giải thuật A* \cite{hart1968astar} bảo đảm tìm được lời giải ngắn
    nhất tối ưu toàn cục, hàm đánh giá heuristic $h(n)$ bắt buộc phải
    thỏa mãn tính chất chấp nhận được ($h(n) \le h^*(n)$) và tính nhất
    quán \cite{dechter1985astar}. Việc xây dựng một hàm heuristic phản
    ánh chính xác khoảng cách Manhattan trong điều kiện khối có thể bao
    phủ hai ô mỗi bước, đồng thời xử lý trơn tru cả giai đoạn khối bị
    phân tách thành hai thực thể độc lập đòi hỏi các phân tích toán học
    nới lỏng bài toán hết sức chặt chẽ \cite{pearl1984heuristics}.
\end{enumerate}

\subsection{Mục tiêu và Kết quả Đạt được của Đồ án}
Đồ án được thực hiện trong khuôn khổ môn học Nhập môn Trí tuệ Nhân tạo
với các mục tiêu trọng tâm sau:
\begin{itemize}[--]
    \item \textbf{Vận dụng lý thuyết tìm kiếm trong không gian trạng thái:}
    Áp dụng các giải thuật tìm kiếm đã học trong môn học  bao gồm cả
    tìm kiếm không thông tin (BFS, DFS, UCS) và tìm kiếm có thông tin
    (A*)  vào việc tự động giải quyết trò chơi giải đố Bloxorz.

    \item \textbf{Mô hình hóa và xây dựng trò chơi hoàn chỉnh:}
    Hiện thực hóa chính xác các quy tắc vật lý và cơ chế tương tác phức tạp
    của Bloxorz: động học lật khối theo 3 tư thế không gian, điều kiện chịu
    tải của gạch cam mong manh, mạng lưới cầu đóng/mở động theo công tắc tròn
    và công tắc chữ X, cùng cơ chế phân thân và tái hợp nhất hai khối lập
    phương độc lập.

    \item \textbf{Thiết kế hàm Heuristic phù hợp:}
    Xây dựng hàm Heuristic cho giải thuật A* dựa trên khoảng cách Manhattan
    kết hợp xử lý sai lệch tư thế và chi phí hợp nhất phân thân, đồng thời
    chứng minh tính chất chấp nhận được ($h(n) \le h^*(n)$) để bảo đảm A*
    luôn tìm ra đường đi ngắn nhất.

    \item \textbf{Đo đạc và đánh giá so sánh thực nghiệm:}
    Thực hiện kiểm thử và thu thập số liệu thực tế trên nhiều màn chơi
    đại diện có độ phức tạp khác nhau; phân tích và đối sánh định lượng các
    chỉ số về thời gian chạy, dung lượng bộ nhớ sử dụng, số trạng thái mở rộng
    và tính tối ưu của lời giải giữa 4 thuật toán.
\end{itemize}

Dựa trên các mục tiêu trên, nhóm đã hoàn thành đầy đủ các nội dung:
\begin{enumerate}[1.]
    \item Xây dựng hoàn chỉnh trò chơi Bloxorz với giao diện trực quan 3D,
    hỗ trợ bản đồ thu nhỏ (Minimap) và bộ nạp màn chơi từ dữ liệu có sẵn.
    \item Cài đặt hoàn chỉnh bộ 4 thuật toán tìm kiếm (BFS, DFS, UCS, A*)
    trên cùng một cấu trúc trạng thái chung, tối ưu hóa việc lưu trữ nút
    để tiết kiệm bộ nhớ trong quá trình duyệt đồ thị.
    \item Chứng minh toán học tính chấp nhận được của hàm Heuristic và kiểm
    chứng tính tối ưu của đường đi thông qua thực nghiệm.
    \item Đo lường, vẽ đồ thị đối sánh và phân tích ưu nhược điểm thực tế
    của từng giải thuật trên tập màn chơi thử nghiệm.
\end{enumerate}